\documentclass[10pt]{article}

\usepackage[utf8]{inputenc}

\usepackage[T1]{fontenc}

\usepackage{amsmath}

\usepackage{amsfonts}

\usepackage{amssymb}

\usepackage[version=4]{mhchem}

\usepackage{stmaryrd}



\begin{document}

степеней свободы. Состояние такой системы обычно характеризуется ее расположением (конфигурацией) и скоростью изменения последнего, а закон движения указывает, с какой скоростью изменяется состояние системы.



В простейших случаях состояние можно охарактеризовать посредством величин $w_{1}, \ldots, w_{m}$, к-рые могут принимать произвольные (действительные) значения, причем двум различным наборам величин $w_{1}, \ldots, w_{m}$ и $w_{1}^{\prime}, \ldots, w_{m}^{\prime}$ отвечают различные состояния, и обратно, а близость всех $w_{i}^{\prime}$ к $w_{i}$ означает близость соответствующих состояний системы. Закон движения тогда записывается в виде автономной системы обыкновенных дифференциальных уравнений:



\[

\dot{w}_{i}=f_{i}\left(w_{1}, \ldots, w_{m}\right), i=1, \ldots, m .

\]



Рассматривая значения $w_{1}, \ldots, w_{m}$ как координаты точки $w$ в $m$-мерном пространстве, можно геометрически представить соответствующее состояние Д.с. посредством этой точки $w$. Последнюю называют фазо в о й (иногда также и зо б р а жа юще й, или пре дставляю щей) то ч кой, а пространство - фазовым пространством системы. (Прилагательное «фазовый» связано с тем, что в прошлом состояния системы нередко называли ее фазами.) Изменение состояния со временем изображается как движение фазовой точки по нек-рой линии (так наз. фазов ой траекто ри и; часто, впрочем, ее называют просто т р а е к т о р и е й) в фазовом пространстве. В последнем определено векторное поле, сопоставляющее каждой точке $w$ выходящий из нее вектор $f(w)$ с компонентами



\[

\left(f_{1}\left(w_{1}, \ldots, w_{m}\right), \ldots, f_{m}\left(w_{1}, \ldots, w_{m}\right)\right) .

\]



Дифференциальные уравнения (1), к-рые с помощью введенных обозначений можно сокращенно записать в виде



\[

\dot{w}=f(w) \text {, }

\]



означают, что в каждый момент времени векторная скорость движения фазовой точки (или, как часто говорят, в е к т о ф фазово й скор ости; не путать с употреблением того же термина в оптике и вообще при рассмотрении различных волновых процессов) равна вектору $f(w)$, исходящему из той точки $w$ фазового пространства, где в данный момент находится движущаяся фазовая точка. В этом состоит так наз. кине матическая инте рпретация системы дифференциальных уравнений (1).



Напр., состояние частицы без внутренних степеней свободы (материальной точки), движущейся в потенциальном поле с потенциалом $U\left(x_{1}, x_{2}, x_{3}\right)$, характеризуется ее положением $x=\left(x_{1}, x_{2}, x_{3}\right)$ и скоростью $\dot{x}$; вместо последней можно использовать импульс $p=m x$, где $m$ - масса частицы. Закон движения можно записать в виде



\[

\dot{x}=\frac{1}{m} p, \quad \dot{p}=-\operatorname{grad} U(x) .

\]



Формулы (4) представляют собой сокращенную запись системы шести обыкновенных дифференциальных уравнений 1-го порядка. Фазовым пространством здесь служит 6-мерное евклидово пространство, шесть компонент вектора фазовой скорости суть компоненты обычной скорости и силы, а проекция фазовой траектории на пространство $x_{i}$ (параллельно пространству импульсов) есть траектория частицы в обычном смысле слова.



В ряде случаев оказывается, что невозможно установить такое соответствие между всеми состояниями Д. с. и точками евклидова пространства, к-рое обладало бы желаемыми свойствами, однако такое соответствие можно установить локально, то есть для состояний, к-рые достаточно близки друг к другу. Сохраняя термин «фазовое пространство» для совокупности всех состояний Д.с., можно сказать, что в общем случае фазовое пространство является не евклидовым пространством, а нек-рым дифференцируемым многообразием $W^{m}$. Локально, то есть в любой карте (локальной системе координат) многообразия $W^{m}$, движение Д. с. описывается системой дифференциальных уравнений вида (1). Глобальное же (то есть пригодное для всех состояний Д. с.) и инвариантное (то есть не зависящее от выбора карты) описание движения дается уравнением (3), в к-ром $f$ является заданным $W^{m}$ векторным полем, сопоставляющим каждой точке $w$ вектор $f(w)$, лежаший в касательном пространстве к многообразию в этой точке; уравнение (3) означает, что в процессе движения фазовая точка, совпадающая в данный момент времени с точкой $w \in W^{m}$, имеет в этот момент скорость $f(w)$. В локальных координатах вектор $f(w)$ представляется посредством своих компонент (2), и (3) сводится к (1).



Даже во многих случаях, когда фазовое пространство является евклидовым, часть движений рассматриваемой Д.с. может описываться посредством векторного поля на нек-ром инвариантном многообразии $W$, то есть подмногообразии фазового пространства, обладающем тем свойством, что траектория, проходящая через какую-нибудь точку $w \in W$, целиком лежит в $W$. Так, уже в предыдушем примере, если ограничиться движениями с определенным значением энергии $E$, нужно рассматривать систему (4) не во всем 6-мерном евклидовом пространстве переменных $(x, p)$, а на его 5-мерном подмногообразии, выделяемом уравнением



\[

\frac{p^{2}}{2 m}+U(x)=E,

\]



где $p^{2}=p_{1}^{2}+p_{2}^{2}+p_{3}^{2}$. Инвариантность этого многообразия отражает тот факт, что при движении частицы в потенциальном поле ее энергия сохраняется, то есть $p^{2} / 2 m+U(x)$ является первым интегралом системы (4) (так наз. и н те гр а л эн е рг и и). Много аналогичных примеров связано с циклическими координатами.



Примером Д. с. с неевклидовым фазовым пространством является твердое тело с неподвижной точкой $O$. Если ввести две ортогональные системы координат с началом в $\mathrm{O}-$ одну неподвижную, а другую жестко связанную с телом, то очевидно, что положение твердого тела будет характеризоваться расположением второй системы координат относительно первой, то есть ортогональной матрицей 3-го порядка с определителем 1 (или каким-нибудь эквивалентным способом). Поэтому совокупность всевозможных положений данной механич. системы (или, как говорят, ее ко нфигурационное пространство) есть специальная ортогональная группа 3-го порядка $S O(3)$. Фазовым же пространством $W^{6}$ является касательное расслоение группы $S O(3)$, ибо скорость изменения положения характеризуется касательным вектором к $S O(3)$. В качестве локальных координат в $S O(3)$ (выбор таковых автоматически определяет и нек-рые локальные координаты в $W^{6}$ ) обычно принимают углы Эйлера; тогда уравнения движения записываются как уравнения Эйлера (движения твердого тела).



Изложенная выше кинематич. интерпретация автономной системы обыкновенных дифференциальных уравнений (1) [или картина движения фазовых точек в фазовом многообразии согласно уравнению (3)] никак не связана с тем, описывают ли эти уравнения какую-либо механич. систему. Поэтому термин «Д. с.» стал применяться в более широком смысле, означая произвольную физич. (скажем, радиотехнич.) систему, описываемую дифференциальными уравнениями вида (1) или (3), а затем - и просто систему дифференциальных уравнений такого вида, безотносительно к ее происхождению. Среди всех Д.с. в этом расширенном смысле слова Д. с. механики выделяются нек-рыми специфич. свойствами; большинство из них относится к специальному классу гамильтоновых систем. (Однако в механике рассматриваются и системы, не входящие в этот класс,таково, напр., большинство неголономных систем. В то же время гамильтоновы системы встречаются и в ряде вопросов физики.)



В этом смысле понятие Д. с. эквивалентно понятию автономной системы дифференциальных уравнений вида (1) или (3). Однако практически о Д.с. говорят тогда, когда





\end{document}